\documentclass[10pt,letterpaper]{article}
\usepackage{cloudsuite}
\renewcommand{\DocShort}{FINOPS / SLA}
\begin{document}
\coverpage{FinOps and SLA Reference Templates}{Metering, forecasting, resilience and contractual assurance}{Concepts \textbullet{} decisions \textbullet{} verification artifacts}
\sourceStatement
\tableofcontents
\newpage
\section{Operational and financial baseline}
This workbench uses the source's Chapter 17 categories (business and usage costs, pricing models) and Chapter 18 categories (availability, reliability, performance, scalability and resiliency). RPO/RTO scenarios and budget thresholds are additional practice fields, not sourced numerical guarantees.
\clearpage\section{F01. Cost model and usage assumptions}
\sourceRef{17}{479}\\[8pt]
\field{Forecast period and service boundary}
\field{One-time migration / tooling / training costs}
\field{Compute quantity and unit rate}
\field{Storage capacity, I/O and retention}
\field{Inbound / outbound / WAN data traffic}
\field{User subscriptions or transaction volume}
\field{Operations / support allocation}
\field{Best, forecast and stress scenario}
\field{Assumptions, rate source and review date}
\begin{workbox}[title={Management review criteria}]
\begin{itemize}
\item Use quantities and units consistently.
\item Account for transfer, support and one-time expense.
\end{itemize}
\end{workbox}
\clearpage\section{F02. Consumption tagging and chargeback}
\sourceRef{17}{479}\\[8pt]
\field{Business owner, application and environment}
\field{Cost center and shared resource allocation rule}
\field{Metering event and collection frequency}
\field{Resource owner labels / account mapping}
\field{Unattributed spend handling}
\field{Budget and overrun threshold}
\field{Finance reconciliation and corrections}
\begin{workbox}[title={Management review criteria}]
\begin{itemize}
\item Reconcile untagged and shared spend.
\item Preserve the raw usage record.
\end{itemize}
\end{workbox}
\clearpage\section{F03. Cloud pricing and commitment review}
\sourceRef{17}{479}\\[8pt]
\field{On-demand baseline and observed utilization}
\field{Reserved / committed allocation scenario}
\field{Term length, cancellation and flexibility constraints}
\field{Workload growth and seasonality}
\field{Break-even utilization and stress case}
\field{Data transfer and exit sensitivity}
\field{Final purchasing approval}
\begin{workbox}[title={Management review criteria}]
\begin{itemize}
\item Do not assume commitment savings without matching usage.
\item Identify obligations after workload retirement.
\end{itemize}
\end{workbox}
\clearpage\section{F04. Service-level objective and agreement}
\sourceRef{18}{503}\\[8pt]
\field{Service and consumer/provider parties}
\field{SLI measurement formula and data source}
\field{SLO threshold and rolling measurement window}
\field{Excluded events / eligible minutes}
\field{Alerting and escalation contacts}
\field{SLA remedy and dispute process}
\field{Reporting retention and evidence owner}
\field{Change review and revalidation date}
\begin{workbox}[title={Management review criteria}]
\begin{itemize}
\item Quality terms are measurable.
\item Exclusion rules cannot be silently modified.
\end{itemize}
\end{workbox}
\clearpage\section{F05. Availability and failure assessment}
\sourceRef{18}{503}\\[8pt]
\field{Eligible observation time (minutes)}
\field{Eligible outage time (minutes)}
\field{Computed availability percentage}
\field{Incident count and failure boundaries}
\field{Detection time and time to switchover}
\field{Mean time between failures (measurement definition)}
\field{Customer impact and credit eligibility}
\begin{workbox}[title={Management review criteria}]
\begin{itemize}
\item Calculations can be reproduced from event timestamps.
\item Incidents and uptime are separately examined.
\end{itemize}
\end{workbox}
\clearpage\section{F06. Continuity restoration scorecard}
\sourceRef{18}{503}\\[8pt]
\field{Recovery scenario and scope}
\field{Backup point timestamp and integrity}
\field{Service interruption start/end}
\field{Observed switchover duration}
\field{Observed total recovery duration}
\field{Observed data-loss age}
\field{Target RPO / RTO and pass/fail}
\field{Reviewer, gap owner and next rehearsal}
\begin{workbox}[title={Management review criteria}]
\begin{itemize}
\item Restore testing validates complete service operation.
\item Time sources and dependency assumptions are recorded.
\end{itemize}
\end{workbox}
\clearpage\section{F07. Monthly service assurance board}
\sourceRef{17}{479}\\[8pt]
\field{Service and business owner}
\field{Service quality measures and month-over-month trend}
\field{Incident and capacity summary}
\field{Measured usage and allocated spend}
\field{Budget variance and anomalies}
\field{Open control exceptions and accountable owners}
\field{Supplier SLA outcomes and follow-up dates}
\begin{workbox}[title={Management review criteria}]
\begin{itemize}
\item Pair spend with service outcomes.
\item Every anomaly has a disposition and owner.
\end{itemize}
\end{workbox}
\section{Worked measurement example}
Suppose an illustrative observation window covers 30 days (43,200 minutes), with 45 eligible outage minutes after the agreed classification. The resulting availability is:
\[
 A=(43{,}200-45)/43{,}200=99.8958\%\quad(\text{approximately}).
\]
These numbers are a hypothetical exercise; they are not an SLA in the source or an external cloud provider guarantee.

\refs
\end{document}
